\documentclass[letterpaper,11pt]{article}

\usepackage{latexsym}
\usepackage[T1]{fontenc}
\usepackage[utf8]{inputenc}
\usepackage{microtype}
\usepackage[empty]{fullpage}
\usepackage{titlesec}
\usepackage{marvosym}
\usepackage[usenames,dvipsnames]{color}
\usepackage{enumitem}
\usepackage[hidelinks]{hyperref}
\usepackage{fancyhdr}
\usepackage[english]{babel}
\usepackage{tabularx}
\usepackage{array}
\usepackage{fontawesome5}
\usepackage{qrcode}
\usepackage{tikz}
\input{glyphtounicode}
\usepackage[sfdefault]{roboto}

\pagestyle{fancy}
\fancyhf{}
\renewcommand{\headrulewidth}{0pt}
\renewcommand{\footrulewidth}{0pt}

\addtolength{\oddsidemargin}{-0.55in}
\addtolength{\evensidemargin}{-0.5in}
\addtolength{\textwidth}{1.1in}
\addtolength{\topmargin}{-.76in}
\addtolength{\textheight}{1.18in}

\urlstyle{same}
\raggedbottom
\raggedright
\setlength{\tabcolsep}{0in}

\titleformat{\section}{\vspace{-2pt}\scshape\raggedright\large\bfseries}{}{0em}{}[\color{black}\titlerule \vspace{-6pt}]
\pdfgentounicode=1

\newcommand{\resumeItem}[1]{\item\small{{#1 \vspace{-1pt}}}}
\newcommand{\resumeSubheading}[4]{
  \vspace{-1pt}\item
    \begin{tabular*}{1.0\textwidth}[t]{l@{\extracolsep{\fill}}r}
      \textbf{#1} & \textbf{\small #2} \\
      \textit{\small#3} & \textit{\small #4} \\
    \end{tabular*}\vspace{-6pt}
}
\newcommand{\resumeProjectHeading}[2]{
    \item
    \begin{tabular*}{1.001\textwidth}{l@{\extracolsep{\fill}}r}
      \small#1 & \textbf{\small #2}\\
    \end{tabular*}\vspace{-6pt}
}
\newcommand{\resumeSubHeadingListStart}{\begin{itemize}[leftmargin=0.0in, label={}]}
\newcommand{\resumeSubHeadingListEnd}{\end{itemize}}
\newcommand{\resumeItemListStart}{\begin{itemize}[leftmargin=0.15in]}
\newcommand{\resumeItemListEnd}{\end{itemize}\vspace{-4pt}}
\renewcommand\labelitemi{$\vcenter{\hbox{\tiny$\bullet$}}$}

\begin{document}

\begin{center}
  {\Huge \scshape Tejas S}\\[-2pt]
  {\small Bengaluru, India}\\[3pt]
  {\small
    \faPhone\ +91-8618904742 \quad
    \href{mailto:tejassureshofficial@gmail.com}{\faEnvelope\ \underline{tejassureshofficial@gmail.com}} \quad
    \href{https://tejassureshofficial.in}{\faGlobe\ \underline{tejassureshofficial.in}} \quad
    \href{https://www.linkedin.com/in/tejas-s-57138816a/}{\faLinkedin\ \underline{LinkedIn}} \quad
    \href{https://github.com/Tez861910}{\faGithub\ \underline{GitHub}}
  }
\end{center}
\vspace{-8pt}

\begin{tikzpicture}[remember picture,overlay]
  \node[anchor=north east] at ([xshift=-0.8cm,yshift=-0.9cm]current page.north east)
    {\qrcode[height=1.4cm]{https://tejassureshofficial.in}};
\end{tikzpicture}

\section{Summary}
\resumeSubHeadingListStart
  \resumeItem{I build product software across web, desktop, and cross-platform applications with a bias toward clarity, maintainability, and real operational usefulness. The common thread in my recent work is software that has to carry actual workflow load: multi-role dashboards, content-heavy company platforms, internal manufacturing tools, and local-first utility suites. I am most useful in product engineering roles where interface decisions, content structure, and implementation details need to reinforce one another.}
\resumeSubHeadingListEnd
\vspace{-10pt}

\section{Product Focus}
\resumeSubHeadingListStart
  \resumeItem{Workflow-oriented interfaces that make dense information easier to use, not just prettier to look at.}
  \resumeItem{Product surfaces that span frontend, backend, and operational rules rather than isolated UI tickets.}
  \resumeItem{Content systems where information architecture matters as much as visuals.}
  \resumeItem{Delivery environments where maintainability, debugging, and future extension are part of the work from day one.}
\resumeSubHeadingListEnd
\vspace{-10pt}

\section{Working Style}
\resumeSubHeadingListStart
  \resumeItem{Understand the domain before optimising the UI so the product solves the actual workflow instead of only looking polished.}
  \resumeItem{Keep frontend structure, backend rules, and content architecture aligned so the software feels coherent to real users.}
  \resumeItem{Prefer maintainable systems with clear extension paths over one-off hero implementations.}
  \resumeItem{Stay comfortable moving between product framing, implementation detail, and final presentation quality.}
\resumeSubHeadingListEnd
\vspace{-10pt}

\section{Core Stack}
\begin{itemize}[leftmargin=0.15in, label={}]
  \small{\item{
    \textbf{Frontend}{: React, Next.js, TypeScript, Vite, Tailwind CSS, Framer Motion} \\
    \textbf{Backend}{: Node.js, Express, Prisma, REST APIs, validation, uploads, security middleware} \\
    \textbf{Desktop}{: .NET 8, C\#, WinUI 3, Win2D, AssimpNet, screenshot and document workflows} \\
    \textbf{Cross-Platform \& Local-First}{: Flutter, Dart, Rust, Flutter Rust Bridge, local-first architecture, multi-platform packaging} \\
    \textbf{Data \& Tooling}{: PostgreSQL, MySQL, SQLite, Git, GitHub Actions, production-oriented repo structure}
  }}
\end{itemize}
\vspace{-12pt}

\section{Experience}
\resumeSubHeadingListStart
  \resumeSubheading
    {Software Engineer --- Product, Web \& Desktop}{Dec 2024 -- Dec 2025}
    {Printalytix}{Bengaluru, India}
    \resumeItemListStart
      \resumeItem{Worked across customer-facing and internal product surfaces, from upload-heavy web flows to operations-focused tooling.}
      \resumeItem{Built and evolved \textbf{MIND} as a .NET 8 Windows application for 3D model inspection, manufacturing guidance, screenshot capture, and report generation.}
      \resumeItem{Delivered React and Node-based systems with protected workflows, validation, uploads, and maintainable product structure.}
      \resumeItem{Contributed across frontend, backend, and workflow design instead of staying in one layer of the stack.}
      \resumeItem{Helped shape software that supports real operational follow-through rather than just presentation.}
    \resumeItemListEnd

  \resumeSubheading
    {Volunteer --- Full-Stack Development}{Oct 2023 -- Jan 2024}
    {Old Dominion University}{Norfolk, VA, USA}
    \resumeItemListStart
      \resumeItem{Built a university threads platform with role-based access, discussion flows, and real-time messaging.}
      \resumeItem{Delivered an examination management system covering assessment workflows, submissions, and administrative coordination.}
      \resumeItem{Worked on backend queries, caching, and application structure to keep the platform responsive as features expanded.}
      \resumeItem{Strengthened early full-stack foundations around auth, validation, role-aware UX, and shared-system design.}
    \resumeItemListEnd
\resumeSubHeadingListEnd
\vspace{-10pt}

\section{Project Highlights}
\resumeSubHeadingListStart
  \resumeProjectHeading
    {\textbf{Pettige} --- Product developer $|$ \emph{2025 -- present}}{}
    \resumeItemListStart
      \resumeItem{Local-first utility suite with \textbf{30+ tools} spanning PDF, image, audio, video, QR, OCR, storage, and developer workflows.}
      \resumeItem{Designed the split architecture between Flutter screens and a Rust core connected through Flutter Rust Bridge.}
      \resumeItem{Built dedicated PDF Studio, Image Studio, Audio Studio, Video Studio, and Dev Workbench flows instead of one generic tool list.}
      \resumeItem{Kept packaging, storage, and OCR paths in scope across Android, Windows, Linux, and Apple-platform targets.}
    \resumeItemListEnd
    \vspace{-8pt}

  \resumeProjectHeading
    {\textbf{Open Solar Toolkit} --- Frontend and product systems developer $|$ \emph{2025 -- present}}{}
    \resumeItemListStart
      \resumeItem{Solar operations platform covering \textbf{7 workstreams} across administration, project management, sales, supply chain, installation, maintenance, and user dashboards.}
      \resumeItem{Approached the product as a set of specialised interfaces inside one shared system rather than a generic one-size-fits-all dashboard.}
      \resumeItem{Built around operational density, with drag-and-drop interaction support, charting, theming, responsiveness, and validation tooling already in scope.}
      \resumeItem{Demonstrates how I structure multi-role products so the shared system stays coherent while each workflow keeps its own shape.}
    \resumeItemListEnd
    \vspace{-8pt}

  \resumeProjectHeading
    {\textbf{MIND Manufacturing Intelligence} --- Desktop product developer $|$ \emph{2025 -- present}}{}
    \resumeItemListStart
      \resumeItem{Windows desktop application for inspecting 3D models and supporting manufacturing decisions in one workflow.}
      \resumeItem{Worked across views, view models, rendering services, assets, screenshot capture, and PDF/reporting paths.}
      \resumeItem{Treated diagnostics and debuggability as first-class concerns because rendering-heavy internal tools need observability to remain maintainable.}
      \resumeItem{Shows range beyond web work into workflow-heavy desktop software where correctness and operator clarity matter.}
    \resumeItemListEnd
    \vspace{-8pt}

  \resumeProjectHeading
    {\textbf{ChemLife Innovations Platform} --- Web platform developer $|$ \emph{2025}}{}
    \resumeItemListStart
      \resumeItem{Content-rich company platform for an animal-nutrition and biotech business with reusable data models for technologies, concepts, awards, and knowledge content.}
      \resumeItem{Structured the site so technical credibility, educational content, and future lead-generation flows can live in one maintainable system.}
      \resumeItem{Treated content architecture as a core engineering concern instead of leaving company information scattered across unrelated pages.}
      \resumeItem{Demonstrates how business websites can become structured product surfaces instead of thin marketing shells.}
    \resumeItemListEnd
    \vspace{-8pt}

  \resumeProjectHeading
    {\textbf{Resolute Solutions Site} --- Frontend developer and content systems builder $|$ \emph{2025}}{}
    \resumeItemListStart
      \resumeItem{Live multi-vertical business site spanning \textbf{4 service lines} with shared navigation, inquiry flow, and distinct presentation per vertical.}
      \resumeItem{Extended the housekeeping side into a more catalogue-like product surface instead of a generic marketing page.}
      \resumeItem{Carried the work through to production delivery with PWA support, build tooling, and deployment considerations.}
      \resumeItem{Shows how one codebase can support multiple visual identities without losing system coherence.}
    \resumeItemListEnd
    \vspace{-8pt}

  \resumeProjectHeading
    {\textbf{Printalytix Platform} --- Full-stack product developer $|$ \emph{2024 -- 2025}}{}
    \resumeItemListStart
      \resumeItem{Full-stack workflow platform for customer intake, uploads, protected operations, and ongoing internal follow-through.}
      \resumeItem{Worked across frontend, backend, and data layers so validation, access rules, and product structure behave like one system.}
      \resumeItem{Helped modernise the frontend workflow while keeping richer form, table, and upload surfaces intact.}
      \resumeItem{Reinforces the kind of work I enjoy most: practical software that supports transactions and operations rather than isolated screens.}
    \resumeItemListEnd
    \vspace{-8pt}

  \resumeProjectHeading
    {\textbf{University Threads Platform} --- Full-stack developer $|$ \emph{2023 -- 2024}}{}
    \resumeItemListStart
      \resumeItem{Earlier multi-user collaboration platform that still shows the foundations of my full-stack approach.}
      \resumeItem{Built around authentication, validation, rate limiting, logging, and role-aware behaviour instead of static page rendering alone.}
      \resumeItem{Useful because it demonstrates the base layer beneath later product work: shared systems, protected routes, and structure.}
    \resumeItemListEnd
    \vspace{-8pt}

  \resumeProjectHeading
    {\textbf{Personal Portfolio \& Resume} --- Designer, writer, and frontend developer $|$ \emph{2025}}{}
    \resumeItemListStart
      \resumeItem{Built the main portfolio as a fast main route with deeper case studies, then added a cockpit mode for an interactive version of the same story.}
      \resumeItem{Structured the project data model so project work can be explained with more context and specificity than a bullet-only resume allows.}
      \resumeItem{Kept the resume workflow in the same repository so the web portfolio and downloadable assets stay aligned.}
    \resumeItemListEnd
\resumeSubHeadingListEnd
\vspace{-8pt}

\section{Education}
\resumeSubHeadingListStart
  \resumeSubheading
    {Bachelor of Computer Applications (BCA) --- Computer Science}{Aug 2018 -- Aug 2024}
    {Bengaluru North University}{Bengaluru, India}
    \resumeItemListStart
      \resumeItem{Built a foundation in DBMS, web development, data structures, and algorithms.}
      \resumeItem{Completed a final-year project centred on full-stack web application development with MySQL.}
    \resumeItemListEnd
\resumeSubHeadingListEnd
\end{document}
